\begin{kbvrule}
\noindent\kbvbullet\hspace{0.5em}ಅ-ಕಾರದಿಂದ ಔ-ಕಾರದವರೆಗಿನ ಸ್ವರಗಳ ಒಟ್ಟು ಸಂಖ್ಯೆ ಹದಿನಾಲ್ಕು. ಇವುಗಳಲ್ಲಿ ಅ, ಇ, ಉ, ಋ, ಎ, ಒಗಳು ಮೊಟಕಾಗಿ (ಅಲ್ಪಕಾಲ) ಉಚ್ಚರಿಸಲು ಬರುವ ಸ್ವರಗಳು. ಈ ಆರು ಸ್ವರಗಳಿಗೆ ಹ್ರಸ್ವ ಸ್ವರಗಳೆಂದು ಹೆಸರು.\par
\noindent\kbvbullet\hspace{0.5em}ಆ, ಈ, ಊ, ೠ, ಏ, ಐ, ಓ, ಔ- ಇವು ನೀಳವಾಗಿ (ದೀರ್ಘವಾಗಿ) ಉಚ್ಚರಿಸಲು ಬರುವ ಸ್ವರಗಳು. ಈ ಎಂಟು ಸ್ವರಗಳಿಗೆ ದೀರ್ಘ ಸ್ವರಗಳೆಂದು ಹೆಸರು.
\end{kbvrule}

\kbvsechead{ವ್ಯಂಜನಗಳು :}

\begin{kbvrule}
\noindent\kbvbullet\hspace{0.5em}ಕ-ಕಾರದಿಂದ ಳ-ಕಾರದವರೆಗಿನ ಅಕ್ಷರಗಳನ್ನು ಸ್ವತಂತ್ರವಾಗಿ ಉಚ್ಚರಿಸಲು ಸಾಧ್ಯವಿಲ್ಲ. ಸ್ವರಗಳು ಸೇರಿದಾಗ ಮಾತ್ರವೇ ಇವುಗಳನ್ನು ಉಚ್ಚರಿಸಲು ಸಾಧ್ಯ. ಈ ಜಾತಿಗೆ ಸೇರಿದ ಅಕ್ಷರಗಳಿಗೆ ವ್ಯಂಜನಗಳೆಂದು ಹೆಸರು. ಇವುಗಳ ಒಟ್ಟು ಸಂಖ್ಯೆ ಮೂವತ್ತನಾಲ್ಕು.\par
\noindent\kbvbullet\hspace{0.5em}ವ್ಯಂಜನಗಳು ಅರ್ಧಾಕ್ಷರಗಳು. ಉದಾಹರಣೆಗೆ, "ಕ" ವ್ಯಂಜನದ ಮೂಲ ರೂಪ "ಕ್"; ಇದರೊಂದಿಗೆ "ಅ" ಸೇರಿದಾಗ "ಕ" ಆಗುವುದು. ಕ್ + ಅ = ಕ. ಇದೇ ರೀತಿಯಾಗಿ ಚ್ + ಆ = ಚಾ, ಟ್ + ಇ = ಟಿ, ತ್ + ಈ = ತೀ, ಪ್ + ಉ = ಪು ಆಗುವುದು.\par
\noindent\kbvbullet\hspace{0.5em}ವ್ಯಂಜನಗಳಲ್ಲಿ ವರ್ಗೀಯ ವ್ಯಂಜನಗಳು ಮತ್ತು ಅವರ್ಗೀಯ ವ್ಯಂಜನಗಳು ಎಂದು ಎರಡು ವಿಧ.\par
\noindent\kbvbullet\hspace{0.5em}ವರ್ಗೀಯ ವ್ಯಂಜನಗಳು : ವ್ಯಂಜನಗಳಲ್ಲಿ ಮೊದಲ ಇಪ್ಪತ್ತೈದು ಅಕ್ಷರಗಳು ವರ್ಗೀಯ ವ್ಯಂಜನಗಳು. ಇವುಗಳನ್ನು ಐದು ವರ್ಗಗಳಾಗಿ ವಿಂಗಡಿಸಲಾಗಿದೆ; ಇಂತಹ ಪ್ರತಿಯೊಂದು ವರ್ಗದಲ್ಲಿರುವ ಅಕ್ಷರಗಳ ಸಂಖ್ಯೆ ಐದು.
\end{kbvrule}

\vspace{5pt plus 4pt minus 1pt}
\begin{center}
{\renewcommand{\arraystretch}{1.24}\setlength{\tabcolsep}{10pt}
\begin{tabular}{l c l}
ಕ ಖ ಗ ಘ ಙ & \ldots\ldots & ಕ-ವರ್ಗ \\
ಚ ಛ ಜ ಝ ಞ & \ldots\ldots & ಚ-ವರ್ಗ \\
ಟ ಠ ಡ ಢ ಣ & \ldots\ldots & ಟ-ವರ್ಗ \\
ತ ಥ ದ ಧ ನ & \ldots\ldots & ತ-ವರ್ಗ \\
ಪ ಫ ಬ ಭ ಮ & \ldots\ldots & ಪ-ವರ್ಗ \\
\end{tabular}}
\end{center}
\vspace{5pt plus 4pt minus 1pt}
